\clearpage
\section{Preference weight: M01, $\beta$ (Group B)}

In V6, $\beta$ is the old-age utility weight in
\[
 (1-\beta)\log c_t^y+
 \frac{\beta}{1-\gamma}\log E_t[(c_{t+1}^o)^{1-\gamma}],
 \qquad A_{US,t}=\beta e_{US,t}.
\]
It is a two-period saving weight. It is neither an annual discount factor nor the
present-bias parameter sometimes also called $\beta$.
For two consumption dates with unnormalized weights $1,d$, normalization gives
$\beta=d/(1+d)$. Setting $d=\delta_a^\Delta$ additionally assumes that one V6
transition represents $\Delta$ years and that endpoint utility appropriately
represents both life stages. Aggregating annual utility within life stages can
require a different bridge.

\begin{evidence}
Hirano, Kishi and Toda, pinned 2025 version \src{HKT25}.
Same two-period utility-weight interpretation; Figure~1 note.
&
$\beta=0.5$: theoretical numerical illustration, not an estimate or empirical
range.
&
Analytical overlapping-generations bubble model and numerical figure.
A model period is a generation; the figure does not estimate its calendar length.
&
Strong definition match, absent empirical discipline. Useful for reproducing
the original mechanism; do not label it a literature-estimated baseline.
\\[5pt]
Gourinchas and Parker (2002) \src{GP02}.
Source $\delta$ discounts annual expected CRRA utility; Table~III, p.~71.
&
Estimated $\delta=0.9598$ (corrected SE $0.0179$); optimal-weight variant
$0.9569$ (SE $0.0150$). These are distinct estimators, not interval endpoints.
&
Simulated-moments lifecycle consumption model with earnings risk, liquidity
constraints and retirement. US CEX 1980--1993, with PSID income dynamics;
annual periods.
&
Conditional endpoint bridge:
$\beta=\delta^\Delta/(1+\delta^\Delta)$.
For $\Delta=30$, these points imply $0.2260$ and $0.2105$ (our calculations).
Different EIS, income timing and household budgets make numerical transfer weak.
\\[5pt]
Laibson et al.\ (2026) \src{LLMRT26}.
Annual discount function; Table~3, frequency in footnote~25.
&
Exponential estimate $\delta=0.9601$ (SE $0.0053$).
Present-biased model: continuation $\delta=0.9891$ (SE $0.0051$) and
present-bias weight $b=0.5305$ (SE $0.1143$).
&
Simulated moments from US SCF 1995--2013; high-school households without
college degrees, with credit cards. Annual lifecycle model; 16 moments for
ages 21--60.
&
The exponential point yields conditional $\beta=0.2277$ at $\Delta=30$.
The present-biased continuation factor is not a clean exponential alternative;
$b\neq\beta_{\mathrm{V6}}$. Do not import either parameter directly.
\\[5pt]
\end{evidence}

\noindent\textbf{Recommendation, distinct from evidence.}
Leave the empirical V6 baseline unassigned until life-stage duration and utility
aggregation are specified. Retain $0.5$ solely as the HKT comparison case.
With the explicitly assumed endpoint bridge and $\Delta=20$--$40$, the two
exponential studies' $\delta=0.9569$--$0.9601$ imply approximately
$\beta=0.1465$--$0.3070$. Adding the present-biased continuation factor as a
separate, qualified scenario extends the upper endpoint to $0.4454$.
These are calculated scenarios, not confidence intervals or evidence that V6
must satisfy them. Mapping confidence: high conditional algebra, low empirical
transferability.

\clearpage
\section{Risk curvature: M02, $\gamma$ (Group A)}

V6's old-age certainty equivalent is
$\{E_t[(c_{t+1}^o)^{1-\gamma}]\}^{1/(1-\gamma)}$.
Thus $\gamma$ controls risk aversion, while the intertemporal elasticity of
substitution (EIS) is fixed at one. Its numerical value is not annualized.
The sufficient-theorem condition $\gamma<1$ is a model restriction, not a
criterion for accepting evidence. Values above one remain relevant candidates;
the displayed formula needs its limiting interpretation at $\gamma=1$.

\begin{evidence}
Gourinchas and Parker (2002) \src{GP02}.
Source CRRA curvature; Tables~III and V, pp.~71,79.
&
Estimated curvature $0.5140$ (SE $0.1707$) or $1.3969$ (SE $0.1137$), depending
on weighting. Specification variants span $0.1278$--$5.2823$; some imply
implausible initial wealth.
&
Annual lifecycle simulated moments; US CEX 1980--1993 and PSID income risks.
Risk curvature and annual discounting are jointly estimated.
&
Risk-curvature interpretation is comparable, but source EIS equals inverse
CRRA, whereas V6 separates risk aversion and fixes EIS at one. Neither the
specification span nor selected low estimates constitute a V6 confidence range.
\\[5pt]
Laibson et al.\ (2026) \src{LLMRT26}.
Source consumption-utility CRRA curvature; Table~3.
&
Estimated $1.4663$ (SE $0.2256$) with exponential discounting;
$1.9355$ (SE $0.4350$) with present bias.
&
Annual US lifecycle model; same SCF population and moments described above.
Curvature is jointly fitted with discount parameters.
&
Supports a low-to-moderate alternative near two. No Epstein--Zin separation;
liquidity constraints and present bias alter effective intertemporal responses.
Do not identify effective EIS mechanically with inverse curvature in the
present-biased model.
\\[5pt]
Barsky et al.\ (1997) \src{BJKS97}, Table~XI, p.~563;
Kimball et al.\ (2008) \src{KSS08}, Table~4.
CRRA over lifetime resources inferred from hypothetical income gambles.
&
BJKS: mean risk aversion $12.1193$, harmonic mean about $4.2$.
KSS, correcting response error: mean $8.2$, median $6.3$, interquartile range
$3.9$--$10.3$, 5th--95th percentiles $1.9$--$20.8$.
&
US Health and Retirement Study, older respondents and spouses.
KSS uses repeated 1992--2002 questions, response-error and status-quo
adjustments. This is lifetime-risk elicitation, not an annual asset-pricing model.
&
The ranges describe population heterogeneity, not estimation confidence.
A representative V6 curvature requires an aggregation argument.
The reciprocal of mean risk tolerance is a harmonic risk-aversion statistic,
not the mean; do not silently use it as representative $\gamma$.
\\[5pt]
Holt and Laury (2002) \src{HL02}.
CRRA utility of experimental cash payoffs; Table~3 and p.~1649.
&
Low-payoff behavior centers roughly around $0.3$--$0.5$; one
choice-implied interval is $0.41$--$0.68$. Risk aversion increases when
real stakes increase. These intervals are not confidence intervals.
&
Incentivized paired-lottery choices by students and faculty at three US
universities; one-shot experiments rather than a calendar-frequency model.
&
Evidence supports retaining subunit-curvature scenarios. Transferability is
weak: small cash-payoff risk, stakes dependence and a selected subject pool
differ from V6 old-age consumption risk.
\\[5pt]
\end{evidence}

\clearpage
\subsection{Asset-pricing estimates, synthesis and proposed risk-curvature scenarios}
\begin{evidence}
Vissing-Jorgensen and Attanasio (2003) \src{VA03}.
Epstein--Zin risk curvature; author-uploaded January~10 extended manuscript,
Table~2A, case~4, p.~14.
&
All-stockholder estimates $8.1,6.3,5.5$ as assumed human-wealth share rises
from $1/3$ to $2/3$ to $0.9$; richest-third estimates $28.4,19.9,16.5$.
Middle all-stockholder case: 95\% interval $[0.3,13.9]$.
&
US CEX 1982--1996 stockholders; semiannual consumption growth with overlapping
monthly observations, asset returns and VAR-based human-wealth adjustments.
Structural asset-pricing estimation.
&
Conceptually closer because risk curvature and EIS are separated. Numerical
mapping remains weak: EIS, wealth aggregation and human-capital assumptions
differ. Numbers are verified in the dated extended manuscript, not asserted
to be an identically numbered table in the shorter published article.
\\[5pt]
Chen, Favilukis and Ludvigson (2013) \src{CFL2013}.
Source $\theta$ is risk aversion; $\rho^{-1}$ is EIS.
Equations~(1)--(2), p.~44; Table~2, p.~62.
&
Stockholder $\theta=20,17$ with 95\% intervals $[0.25,40]$, $[1,43.3]$;
aggregate $\theta=57.5,60$ with intervals $[27.5,129]$, $[42,144]$.
EIS estimates exceed one.
&
Quarterly US data 1952Q1--2005Q1; consumption, Treasury bills and equity
portfolios. Profile sieve minimum distance with second-step GMM.
Stockholder consumption is a fitted mimicking factor.
&
Map $\gamma=\theta$ conceptually, without annualization. Infinite-horizon
continuation-value risk and freely estimated EIS differ from V6.
Wide bootstrap intervals and possible model misspecification limit transfer.
The span $17$--$60$ is between specifications, not a confidence interval.
\\[5pt]
Elminejad, Havranek and Irsova (2025) \src{EHI25}.
Relative-risk-aversion meta-analysis; Table~6 and main discussion.
&
Bias-adjusted summaries differ by context: roughly one in economics,
two--seven in finance. US-context prediction $5.20$, 95\% interval
$[2.28,8.11]$.
&
Synthesis of 1,021 estimates from 92 studies, with literature search ending
May~2022. Studies differ in populations, frequencies and empirical designs.
This is secondary synthesis, not a new original estimate.
&
Useful context for choosing a sensitivity grid; not a direct V6 measurement.
Its discipline split cautions against pooling lifecycle and asset-pricing
evidence into a single consensus number.
\\[5pt]
\end{evidence}

\noindent\textbf{Recommendation, distinct from evidence.}
Use $\gamma=5$ as a provisional finance-oriented candidate, with $2$ as an
explicit lifecycle-oriented alternative. A core sensitivity grid
$\{0.5,2,5,6.3,10\}$ preserves low experimental and higher survey/finance
evidence; add a separate high-risk family $\{17,20,60\}$ from recursive
asset-pricing estimates. This grid is a research recommendation, not a pooled
credible interval. Conceptual mapping confidence is medium; numerical
confidence is low. The candidate $5$ violates the sufficient-theorem
restriction $\gamma<1$; retain and report that conflict instead of choosing
a low number solely to preserve the theorem.
